\section{Introduction}
\label{sec:intro}

Many data management tasks rely on \emph{data quality constraints}, which define what constitutes valid data.
However, such constraints are often missing, incomplete, or outdated in real-world datasets. This gap has motivated a large body of work on \emph{data profiling}, where data constraints are automatically discovered from the data itself~\cite{profSurveyFelix2015}. 
Profiling algorithms extract constraints, such as functional dependencies (FDs) and denial constraints (DCs), by identifying patterns that hold over a given data instance, and have been widely used to support data cleaning and error detection~\cite{dependencyTree2020}.

However, a fundamental assumption underlying this line of work is that constraints must be inferred from data, and thus, they are limited to discovering constraints that reflect what the data \emph{happens to satisfy}, not what it is \emph{required to satisfy}, given some policy, regulation, or legislation. No profiling algorithm can make that distinction, nor can it surface constraints that must hold by design but are violated in the data.
\altic{[R3-O1] Suavizar esta afirmação: ela foi lida como factualmente incorreta. (i) Existe muito trabalho guiando especialistas a definir/refinar regras, e (ii) a descoberta \emph{aproximada} de DCs (``Discovery of Approximate (and Exact) Denial Constraints'', VLDB 2019) já descobre regras que NÃO valem plenamente nos dados. Reescrever para distinguir ``regras requeridas por norma'' de ``regras que valem (exata ou aproximadamente) nos dados'', e citar a descoberta aproximada de DCs aqui.}

Many datasets are governed by rules specified in regulatory documents, policy manuals, and operational specifications, which define requirements such as mandatory fields, valid values, eligibility criteria, uniqueness conditions, quotas, and ordering rules~\cite{batini2009methodologies,serra2024IQcontextDtaQuality,CABALLERO2022102058}. Formalizing these rules as constraints is difficult because they depend on combinations of attributes, relationships across tuples, comparison operators, thresholds, and context-dependent conditions, requiring expressive \emph{data quality constraints}. However, since such rules are typically expressed only in natural language, they are rarely translated into executable constraints and used with the data they govern. Consequently, datasets are frequently released without an explicit representation of their own correctness requirements, so violations remain undetected and subsequent analyses may be based on data that does not satisfy their intended specification.

This problem is common in modern data settings, where open-data portals publish standalone files with limited metadata, data lakes aggregate heterogeneous sources, and data pipelines connect systems whose semantics are only partially documented~\cite{datasetSearch2020,Governor2023}.


\begin{example}[Brazilian public university admission]
\label{ex:intro}
Brazil's ``Sistema de Sele\c{c}\~{a}o Unificada'' (SISU) manages the public university placements for millions of students annually. The application data is publicly available,\footnote{\url{https://dadosabertos.mec.gov.br/sisu}} with sensitive fields anonymized. The dataset schema includes attributes such as \emph{student\_id}, \emph{course\_code}, \emph{shift}, \emph{institution\_code}, \emph{year}, and \emph{edition}. The rules that define correct enrollment outcomes, however, are described in separate documents that do not accompany the dataset. For instance, Normative Ruling No.\ 21/2012\footnote{\url{https://sisugestao.mec.gov.br/docs/portaria-2012-21.pdf}} is a PDF file found on the Brazilian government website that states (Article 19):\footnote{Schema elements and quoted statements translated from Portuguese for readability.}

\begin{quote}
\textit{``$\ldots$ the student will be ranked in descending order of scores for the seat option they applied to, respecting the available seats at the institution, by campus location, course, and shift, as well as the student's admission category..''}
\end{quote}

This passage defines an \emph{ordering constraint}: within each context defined by institution, campus, course, shift, and admission category (and per year/edition), a higher score should not result in a worse rank. The following SQL query checks for violations of this requirement in SISU's application data:

\begin{lstlisting}[language=SQL, frame=none, basicstyle=\footnotesize\ttfamily, mathescape]
SELECT t.student_id, t$^\prime$.student_id
  FROM applications t, applications t$^\prime$
 WHERE t.student_id <> t$^\prime$.student_id
   -- same context
   AND t.year = t$^\prime$.year
   AND t.edition = t$^\prime$.edition
   AND t.institution_id = t$^\prime$.institution_id
   AND t.campus_id = t$^\prime$.campus_id
   AND t.course_code = t$^\prime$.course_code
   AND t.shift = t$^\prime$.shift
   AND t.application_mode = t$^\prime$.application_mode
   -- violation: higher score but worse ranking
   AND t.score > t$^\prime$.score
   AND t.ranking > t$^\prime$.ranking;
\end{lstlisting}

Any nonempty result set for this query indicates a violation of a legal requirement. \DIFaddbegin \DIFadd{On the released data, the query returns 76 violating pairs involving 61 applications.
}\DIFaddend \altic{[R2-O1] Reportar AQUI o número concreto de tuplas do SISU que violam o Artigo 19 (esta query). Toda a motivação assume que datasets publicados violam silenciosamente suas normas; dar o número empírico do exemplo motivador torna a motivação crível. O número existe na Tabela~\ref{tab:datasets} (487 violações no total); extrair o valor específico desta regra.}

\qed
\end{example}




In this paper, we introduce the \emph{Text-to-Constraints} problem: given a dataset and a set of normative requirements derived from related regulatory text, \subs{extract data-quality constraints grounded in that text to validate the data}{map those requirements to executable data-quality constraints over the dataset schema}. \comm{ALTI, 2026-08-26: Reframed the problem around grounding given normative requirements; extraction of requirements from regulatory text is outside the current scope.} This problem differs from traditional constraint discovery in that, instead of inferring constraints from patterns observed in the data instance, which may reflect incidental regularities rather than intended rules, it requires reasoning over regulatory semantics, schema structure, and violation logic under a chosen constraint formalism. The \emph{Text-to-Constraints} problem is related to the well-known Text-to-SQL~\cite{llmTEXT2SQLSurvey2025} problem as both involve mapping natural-language specifications to structured representations grounded in a database schema. However, Text-to-SQL generates queries to answer specific user questions,  whereas Text-to-Constraints \subs{extracts declarative invariants}{maps normative requirements to executable constraints} that must hold for all valid dataset instances. \comm{ALTI, 2026-08-26: Aligned the Text-to-SQL comparison with the revised problem definition.} %, yielding executable validation logic with well-defined violation semantics.

\smallskip
\noindent\textbf{Challenges.}
\rem{Regulatory documents are typically long and heterogeneous, and only a fraction of their sentences carry normative force~\cite{https://doi.org/10.1111/jels.12364} (i.e., sentences that prescribe what the data is required to satisfy). In SISU, for example, the normative text spans several pages across multiple articles, yet only a handful of passages describe conditions that can be evaluated against the released data. The remaining passages consist of section titles, definitions, procedural guidance, rationale, and contextual descriptions, which are essential for human readers but do not directly imply data validity. Rules are interleaved throughout this material, often qualified by scope conditions, numeric thresholds, or exception clauses, and expressed in the varied, imprecise language that characterizes legal and administrative text. Translating such passages into formal constraints requires identifying which sentences carry normative force by resolving implicit references to schema elements.}
\add{\subs{Regulatory documents are formed by natural language sentences that carry normative force}{Regulatory documents consist of natural-language text, including sentences that carry normative force}~\cite{https://doi.org/10.1111/jels.12364} (i.e., sentences that prescribe what the data is required to satisfy). Thus, the first challenge we address is mapping these natural-language rules that define valid data to executable data-quality constraints. Example~\ref{ex:intro} illustrates this challenge: Article~19 of the SISU regulation states an ordering requirement in prose, while checking it requires making explicit the relevant context, attributes, tuple comparisons, and violation condition. More generally, a normative passage may combine conditions, qualifiers, cross-record relationships, numeric thresholds, and exceptions, while leaving implicit how its terms correspond to elements of the dataset schema. Mapping such a passage to an executable constraint therefore requires resolving its semantics, identifying the relevant schema elements, and specifying what constitutes a violation.}
\comm{ALTI, 2026-08-26: Reframed the first challenge around the text-to-constraint task and retained SISU as a concrete example, without claiming to identify normative passages in an entire regulatory document.}
%Translating these passages into machine-readable data quality checks is difficult because regulatory text is written for human readers, not for automated processing. Most of the text consists of descriptions, section titles, definitions, rationale, and procedural guidance, with rules interleaved throughout, often qualified by scope conditions, thresholds, or exceptions expressed in varied language. 
%, and difficult to identify without extensive manual annotation.

A second challenge is determining the correct scope of each constraint. 
Omitting qualifiers results in overly broad constraints, flagging valid data as violations.
Adding conditions to the rule results in constraints that are too narrow, missing violations. In Example~\ref{ex:intro}, omitting qualifiers such as year/edition or admission category mis-scopes the check and can flag valid outcomes. On the other hand, if one treats ``respecting the limit of available seats'' as \texttt{AND t.ranking <= seats AND t$^\prime$.ranking <= seats} and only compares applicants within the admitted prefix, then scores involving waitlisted applicants are never checked.

\altic{[R1-O2, R3-O6] Reescrever este ``terceiro desafio'': foi lido como confuso por dois revisores. Problema: ``SQL é um moving target'' não ficou claro, e como toda DC se escreve em SQL, parece contradizer a escolha de DCs como saída. Esclarecer que o ponto NÃO é SQL-vs-não-SQL, mas escolher um espaço-alvo canônico/restrito (DCs) em vez do espaço aberto de formulações SQL equivalentes. Ou seja, a ``moving target'' é a falta de forma canônica do SQL para extração/comparação. Conectar explicitamente com o argumento ``Why DCs'' das Preliminaries.}
A third challenge is the output formalism, and there is an unavoidable compromise. General SQL is expressive, but it is a moving target for extraction, making it difficult to standardize checks across documents. More restricted languages are easier to capture and evaluate, but they miss many cross-record rules, including ordering, quotas, and conditional uniqueness. Example~\ref{ex:intro} requires comparing applicants within the same context, so we need a constraint language that supports such comparisons and can still be checked efficiently on large datasets.


\subs{Generally speaking, the Text-to-Constraints problem can be seen as a form of information extraction, in which systems have traditionally required extensive domain-specific annotation and feature engineering~\cite{ieSurvey2024}.}{Generally speaking, the Text-to-Constraints problem is to map natural-language regulatory text to executable data-quality constraints grounded in a target dataset schema.} \add{This mapping requires linking the terms of a normative passage to schema elements, expressing its semantics in a constraint formalism, and specifying executable violation conditions.} \comm{ALTI, 2026-08-26: Replaced the information-extraction framing with a direct description of the task addressed by Norma.} \subs{Large language models have reduced this burden in related settings, such as checking whether software requirements comply with regulations~\cite{compliance2025ACMTSEM}.}{Large language models provide a promising way to perform this mapping from natural-language regulatory text to structured constraints, as demonstrated in related regulatory compliance tasks~\cite{compliance2025ACMTSEM}.}  
But Text-to-Constraints requires more than identifying relevant passages. The output must be linked to the dataset schema, conform to a constraint formalism, and express violation conditions that can be executed against data. None of these requirements follows automatically from language model outputs. They require explicit mechanisms for schema alignment, predicate grounding, and violation semantics.
\altic{[R1-O1] Posicionar a intro mais a fundo frente ao estado-da-arte, com citações: (i) regras declarativas tipo REEs (Rule/Entity Enhancing rules) que usam ML-predicates, relacionando com nosso espaço de predicados; (ii) como métodos de descoberta de DC poderiam ser estendidos vs. nossa abordagem baseada em texto. Hoje a intro só contrasta com Text-to-SQL e IE. Mover/citar material do Related Work para cá.}


%Text-to-Constraints, however, demands more than identifying relevant passages. The output must be linked to the dataset schema, conform to a constraint formalism, and define violation conditions that support execution. These requirements do not follow automatically from language model outputs. 
%Text-to-Constraints, however, demands more than extracting relevant text passages. The output must conform to a constraint formalism, map to the target schema, and carry well-defined violation semantics, properties that do not follow automatically from language model outputs.
%%% rescrito abaixo invertendo as mesmas frases

\smallskip
\noindent\textbf{From text to data quality constraints.}
These challenges show that Text-to-Constraints goes beyond simple text extraction. A useful solution must identify passages with normative requirements, link them to the dataset schema, and translate them into formal constraints with the correct scope and violation semantics.

To address these challenges, we propose \system, the first framework for Text-to-Constraints. Given a regulatory document and a target dataset, \system\ identifies sentences with normative force and maps them to formal constraints over the dataset schema. \system\ uses denial constraints to represent constraints, as they provide a striking balance between expressiveness and control. Unlike Text-to-SQL systems, which generate ad hoc SQL queries in response to user questions, \system\ derives reusable constraints with well-defined violation semantics. It also differs from traditional data profiling. In data profiling, denial constraints are defined over a narrower set of predicates for tractability. Because \system\ operates over text instead of searching for regularities in the data, it can use a richer predicate vocabulary.\altic{[R3-O5] Termo ``regularities'' foi apontado como vago. Trocar por algo preciso, p.ex. ``padrões estatísticos que valem na instância de dados'' (em contraste com requisitos normativos), deixando claro o contraste dados-vs-texto.} This includes multiple operators, column relationships, null checks, domain checks, and pattern predicates. As a result, \system\ can capture a broader range of regulatory rules than traditional profiling methods, while still keeping the output space structured and restricted.

\system\ supports multiple extraction and mapping modes with different computational-cost trade-offs, ranging from direct, prompt-based variants to structured, multi-step pipelines. Our experiments show that this framework recovers rules that state-of-the-art data profiling algorithms and Text-to-SQL systems do not capture. \system\ identifies real compliance failures in publicly released data using constraints extracted from governing text. A key advantage of this setting is that each detected violation can be traced back to the source passage that justifies the constraint.


A summary of our contributions are the following:
% \smallskip
% \noindent\textbf{Contributions.}
\begin{itemize}[leftmargin=0.35cm]
\item A formalization of the \emph{Text-to-Constraints} problem;
\item \system, the first framework for this problem, with sentence-level \DIFdelbegin \DIFdel{and document-level extraction }\DIFdelend \DIFaddbegin \DIFadd{classification and mapping }\DIFaddend implementations with different cost-effectiveness trade-offs\DIFaddbegin \DIFadd{, and a compilation of the extracted constraints into SQL violation queries}\DIFaddend ;
\item A set of datasets with annotated constraints, and evaluation measures adequate for the Text-to-Constraint problem.
\item An extensive empirical study across multiple models, prompting strategies, and baselines from data profiling and Text-to-SQL, showing that \system\ recovers critical constraints missed by data-only profiling algorithms.
\end{itemize}

The remainder of this paper is organized as follows. Section~\ref{sec:problem} formalizes the Text-to-Constraints problem. Section~\ref{sec:overview} introduces \system, whose classifier and mapper components are detailed in Sections~\ref{sec:classifier} and~\ref{sec:mapper}. Section~\ref{sec:experiments} presents our experimental evaluation of \system\ against different baselines. Section~\ref{sec:related} discusses related work, and Section~\ref{sec:conclusion} presents our conclusions.








%%old stuff

%Integrity constraints, such as keys and domain checks, are usually defined during schema design and enforced by most database management systems (DBMSs). However, many critical data rules involve multiple attributes, records, and comparison operators. These cannot be captured by the common DBMS constraints alone and require richer languages to define \emph{data quality constraints}.
%
%Most datasets found ``in the wild'' already lack traditional DBMS-friendly integrity constraints, let alone business rules that should be captured by non-trivial data quality constraints. Data lakes aggregate information from diverse sources with little oversight, open-data portals publish datasets with minimal metadata, and data pipelines connect rapidly changing systems. In those environments, specifying data quality constraints is a constant challenge~\cite{datasetSearch2020,Governor2023}.
%
% Data quality criteria exist for, if not all, most datasets, but in practice, they are rarely explicitly articulated or defined as well-defined constraints. For instance, many regulatory documents, policy manuals, and operational specifications related to a given dataset define what valid data is for that dataset at some level. Many of these documents cover a variety of data requirements such as mandatory fields, valid ranges, eligibility, uniqueness, quotas, and ranking rules~\cite{batini2009methodologies,serra2024IQcontextDtaQuality,CABALLERO2022102058}. Yet these rules are written in natural language and almost never translated into formal constraints. Datasets are released as standalone files while their governing rules remain in separate documents, so inconsistencies go undetected and analyses rely on data that violates its own rules.

%Many data management tasks rely on \emph{data quality constraints} that describe what valid data looks like. 

%Many datasets are governed by rules that are not captured in their schema. Although most data management systems can enforce basic integrity constraints, such as keys and domain checks, many requirements for determining whether data is valid go beyond what such constraints can express. They may depend on combinations of attributes, relationships across tuples, comparison operators, thresholds, and context-specific conditions. Such requirements call for richer \emph{data quality constraints}. In practice, however, these constraints are often not encoded in the database itself, but instead described in external documents associated with the data.

%This is common in modern data settings. Open-data portals publish standalone files with little metadata, data lakes aggregate heterogeneous sources, and data pipelines connect systems whose semantics are only partially documented~\cite{datasetSearch2020,Governor2023}. Yet the criteria that define valid data often do exist: in regulatory documents, policy manuals, and operational specifications that state mandatory fields, valid values, eligibility requirements, uniqueness conditions, quotas, and ordering rules~\cite{batini2009methodologies,serra2024IQcontextDtaQuality,CABALLERO2022102058}. The problem is that these rules are written in natural language and are rarely translated into formal constraints. As a result, datasets are released without the constraints that govern them, violations remain undetected, and analyses may rely on data that does not satisfy its own governing rules.

%% IF we test against dat acleaning use this
%Current approaches for improving data quality generally do not use texts as a source for grounding constraints. Instead, existing work typically follows one of two routes.
%
%The first is \emph{data profiling}, where algorithms mine constraints within a given formalism, such as functional dependencies (FDs) and denial constraints (DCs) from the data instance~\cite{profSurveyFelix2015}. Profiling is effective at surfacing data patterns, but it cannot distinguish intended rules from coincidental patterns and often returns large numbers of spurious constraints~\cite{martin2025false}. In addition, profiling cannot recover constraints that hold by design but are violated in the released data, whereas regulatory grounding enables the detection of precisely such discrepancies.
%The second is \emph{data cleaning}, where systems detect and repair errors using explicit constraints, learned statistical signals, or both~\cite{dataerrorsZiawash}. Constraint-based cleaning is auditable since flagged records and repairs can be traced to specific constraint violations, but it requires experts to author rules upfront~\cite{chu2013holistic}. Data-driven cleaning reduces manual rule specification, but it can misclassify rare but valid values as errors, and its effectiveness can vary sharply across error types and rates~\cite{cleaningDeploy2024}.
%In both cases, correctness is either inferred from the data instance or assumed to be encoded \emph{a priori} as constraints.
%Beyond these conceptual limitations, most profiling and cleaning systems do not scale to the dozens of columns typical of real-world datasets, further limiting their practical applicability.

%When constraints are unavailable, the typical approach is to discover them from data using \emph{data profiling}~\cite{profSurveyFelix2015}. Profiling algorithms discover constraints within a given formalism, such as functional dependencies (FDs) and denial constraints (DCs), by identifying patterns that hold over the data instance. The fundamental limitation is that a discovered constraint reflects what the data \emph{happens to satisfy}, not what a legislator or policy \emph{requires} it to satisfy. No profiling algorithm can make that distinction, nor can it surface constraints that hold by design but are violated in the data.%.---precisely where regulatory grounding matters most.

%
%
%
%\smallskip
%\noindent\textbf{Limitations of data profiling.}
%Profiling algorithms discover constraints directly from data, targeting formalisms such as functional dependencies, inclusion dependencies, and denial constraints~\cite{profSurveyFelix2015}. While these algorithms have become faster, their outputs are often dominated by spurious constraints that hold due to coincidental patterns or overfitting~\cite{martin2025false}. Because profiling is purely data-driven and ignores external knowledge, it cannot reliably separate intended rules from accidental regularities.
%
%\smallskip
%\noindent\textbf{Limitations of data cleaning.}
%Cleaning systems aim to detect and repair data errors~\cite{dataerrorsZiawash}. Constraint-driven methods are interpretable--each flagged error or repair can be traced to an explicit rule violation--but they require experts to specify rules upfront~\cite{chu2013holistic}. Data-driven methods can reduce the need for constraints, but they often struggle to distinguish rare-but-valid values from true errors~\cite{peng2022garf,raha2019}. Despite substantial progress, deployment remains challenging at scale: large-scale evaluations highlight pitfalls in evaluation metrics, strong dependence on error types and rates, and limited guidance for selecting repair methods; in many settings, automatic repair can even increase the number of errors in the data~\cite{cleaningDeploy2024}.
